\chapter{Marco teórico}\label{cap:marco}

\section{La ciudad como sistema estocástico}

La economía urbana clásica explica la forma de la ciudad como un equilibrio entre el costo de
transporte y el precio del suelo. En el modelo monocéntrico de \citet{alonso1964}, \citet{mills1967} y
\citet{muth1969}, todos los empleos están en un centro y cada hogar elige su distancia a él
compensando suelo más barato con viajes más largos; la renta del suelo decrece con la distancia y la
densidad se ajusta a ella. Dos ideas de esa tradición sostienen este trabajo. La primera es que la red
de transporte y la localización de las actividades se determinan conjuntamente: un corredor rápido
amplía el área accesible, y una política que acerque empleo y vivienda reduce la necesidad de
viajar. La segunda es que la ciudad se construye bajo incertidumbre y con inversiones irreversibles.

\citet{capozza1990} formalizaron esa segunda idea en \emph{La ciudad estocástica}: cuando la renta
urbana sigue un movimiento browniano, la conversión de suelo rural en urbano es una opción, y su
ejercicio óptimo exige que la renta supere el costo de conversión por una prima que crece con la
volatilidad. \citet{capozza1994} extendieron el argumento a la estructura de riesgo de los mercados de
suelo. Por otro lado, la ley de Gibrat de crecimiento proporcional aleatorio, que en tiempo continuo
es exactamente un MBG, es la explicación canónica de la ley de Zipf para el tamaño de las ciudades
\citep{gabaix1999}. Modelar la demanda de una zona urbana como un MBG tiene, por tanto, fundamento en la
literatura de economía urbana y no solo en la de finanzas.

\section{El movimiento browniano geométrico}

\begin{definition}[MBG multivariado]
Sea $W=(W_1,\dots,W_m)$ un movimiento browniano estándar en $\R^m$ y $L\in\R^{Z\times m}$. El proceso
$X=(X_1,\dots,X_Z)$ es un MBG multivariado si
\begin{equation}\label{eq:mbg}
\frac{dX_i(t)}{X_i(t)}=\mu_i\,dt+\sum_{k=1}^{m}L_{ik}\,dW_k(t),\qquad X_i(0)>0,
\end{equation}
con matriz de covarianza instantánea $\Sigma=LL^\top$.
\end{definition}

Por el lema de Itô aplicado a $\ln X_i$ \citep{oksendal2003}, la solución es
\begin{equation}\label{eq:mbg_sol}
X_i(t)=X_i(0)\exp\!\Big(\big(\mu_i-\tfrac12\Sigma_{ii}\big)t+\textstyle\sum_k L_{ik}W_k(t)\Big),
\end{equation}
de modo que $\ln X(t)$ es normal multivariado con media $\ln X(0)+(\mu-\tfrac12\diag\Sigma)t$ y
covarianza $\Sigma t$. En particular $\E[X_i(t)]=X_i(0)e^{\mu_i t}$ y
$\operatorname{Var}[\ln X_i(t)]=\Sigma_{ii}t$. Tres propiedades hacen del MBG un punto de partida
natural para la demanda: es positivo, sus incrementos relativos son independientes del nivel (ley de
Gibrat) y su dispersión crece con el horizonte, lo que traduce en números la intuición de que el
futuro lejano es más incierto que el cercano.

Sus limitaciones también son conocidas y se discuten en el \cref{cap:discusion}: el MBG no tiene
saturación ni reversión a la media, y no captura choques discretos. Para lo último se usa la extensión
de \citet{merton1976}, que añade un proceso de Poisson compuesto:
\begin{equation}\label{eq:merton}
\frac{dX_i}{X_{i^-}}=\big(\mu_i-\lambda\kappa\big)dt+\sum_k L_{ik}\,dW_k+\big(e^{J}-1\big)dN,
\qquad J\sim\mathcal N(m_J,s_J^2),\;\kappa=e^{m_J+s_J^2/2}-1,
\end{equation}
donde la corrección $-\lambda\kappa$ (\emph{compensación}) mantiene la media del proceso igual a la del
MBG sin saltos.

\section{Simulación de Monte Carlo}

Si $g$ es el costo social de un diseño como función de la trayectoria de demanda, el objetivo es
$\E[g(X)]$ y otras funcionales de la distribución de $g(X)$. El estimador de Monte Carlo
$\hat\mu_n=\frac1n\sum_{r=1}^n g(X^{(r)})$ es insesgado y, por el teorema central del límite, su error
es del orden $\sigma_g/\sqrt n$ \citep{glasserman2004}. Dos técnicas reducen la constante o mejoran la
tasa:
\begin{itemize}[nosep]
\item \emph{Variables antitéticas}: se usan los pares $(Z,-Z)$ de normales; si $g$ es monótona en $Z$,
      la covarianza negativa entre los pares reduce la varianza del promedio.
\item \emph{Cuasi-Monte Carlo}: se reemplazan los uniformes pseudoaleatorios por una secuencia de baja
      discrepancia (Sobol) aleatorizada; para integrandos suaves, la varianza del estimador decrece más
      rápido que $1/n$ \citep{owen1997}.
\end{itemize}
Cuando se comparan diseños, la técnica decisiva es la de \emph{números aleatorios comunes}: todos los
diseños se evalúan con los mismos escenarios, de modo que la diferencia de costos entre dos diseños
tiene una varianza mucho menor que la suma de sus varianzas.

\section{Medidas de riesgo}

Sea $L$ una pérdida (aquí, el costo social en valor presente). El valor en riesgo al nivel $\alpha$ es
el cuantil $\VaR_\alpha(L)=\inf\{c:\Prob(L\le c)\ge\alpha\}$, y el valor en riesgo condicional es la
media de la cola,
\begin{equation}\label{eq:cvar}
\CVaR_\alpha(L)=\min_{c\in\R}\Big\{c+\frac{\E[(L-c)^+]}{1-\alpha}\Big\},
\end{equation}
representación de \citet{rockafellar2000} cuyo minimizador es el VaR. A diferencia del VaR, el CVaR es
una medida coherente (en particular, subaditiva) y convexa, y su versión muestral entra de forma lineal
en problemas de optimización. Con $n$ escenarios, el estimador usado es
$\widehat{\VaR}+\frac{1}{n(1-\alpha)}\sum_r(L_r-\widehat{\VaR})^+$.

\section{Opciones reales}

Una inversión irreversible cuyo valor fluctúa es análoga a una opción de compra americana: invertir
equivale a ejercer y destruye la opción de esperar. Para un proyecto de valor $V$ que sigue un MBG con
deriva $\mu$ y volatilidad $\sigma$, costo $I$ y tasa de descuento $r>\mu$, \citet{mcdonald1986}
mostraron que el valor de la oportunidad de invertir es
\begin{equation}\label{eq:ms}
F(V)=\begin{cases}(V^*-I)\,(V/V^*)^{\beta}, & V<V^*,\\ V-I, & V\ge V^*,\end{cases}
\qquad V^*=\frac{\beta}{\beta-1}I,
\end{equation}
con $\beta>1$ la raíz positiva de $\tfrac12\sigma^2\beta(\beta-1)+\mu\beta-r=0$. La regla clásica del
VPN («invertir si $V>I$») es subóptima: el umbral óptimo $V^*$ supera a $I$ en un múltiplo que crece con
la volatilidad \citep{dixit1994}. En transporte, este marco se ha usado para el momento de ampliar
carreteras bajo incertidumbre de demanda \citep{wang2017} y para la flexibilidad frente al cambio
climático \citep{li2022}.

Cuando el valor del proyecto depende de una red y no tiene forma cerrada, la opción se valora por
simulación. El algoritmo de mínimos cuadrados de Monte Carlo (LSM) de \citet{longstaff2001} recorre
hacia atrás las fechas de decisión y estima el valor de continuar mediante una regresión del valor
futuro realizado sobre funciones de la variable de estado; se ejerce cuando el valor inmediato supera
la continuación estimada.

\section{Equilibrio de transporte}

La forma de cuatro etapas (generación, distribución, partición modal y asignación) sigue siendo el
esqueleto de los modelos de demanda \citep{ortuzar2011}. La distribución se modela aquí con un modelo
gravitacional doblemente acotado, la partición modal con un logit binario y la congestión vial con la
función del \citet{bpr1964},
\begin{equation}\label{eq:bpr}
t(v)=t_0\Big(1+\alpha\big(v/c\big)^{\beta}\Big),
\end{equation}
con $\alpha=0{,}15$ y $\beta=4$. El equilibrio entre la elección de modo y los tiempos que esa elección
produce es un problema de punto fijo; cuando el costo de cada modo depende solo de su propia carga y las
funciones son crecientes, equivale a un programa convexo con solución única \citep{sheffi1985}, y se
resuelve con promedios sucesivos, aquí con paso autorregulado \citep{liu2009}.
